\documentclass{nfr}

% This LaTeX template is based on the official Word template found here: https://www.forskningsradet.no/siteassets/utlysninger/vedlegg-utlysninger/template-for-project-description-researcher-project-for-experienced-scientists-and-early-career-scientists.docx
% See the README.md file for more information

% Here's an example bibliography using the numeric (Nature) style and two
% columns for the reference list.

% \usepackage{multicol}
% \renewcommand{\bibpreamble}{\begin{multicols*}{2}}
% \renewcommand{\bibpostamble}{\end{multicols*}}

% \usepackage{biblatex}
% \addbibresource{references1.bib}

% \citestyle{nature}
% \citestyle{}
\usepackage{comment}

% \bibliographystyle{bst/veryshortnum}
\bibliographystyle{apalike}

\title{Application for Faculty Position in Applied Economics}
\author{Zhongli Wang}

\begin{document}

% Project description template
% Researcher Project for Experienced Scientists/Early Career Scientists

% This template consists of
% 
% \textbf{A:} Overview of the chapters and sections to be used as the
% structure of the project description
% 
% \textbf{B:} General guidance
% 
% \textbf{C:} Guidance on the content in each chapter and section
% 
% \textbf{A: Chapters and sections to be used}

% I used pandoc to translate the word template to tex. The formatting of the
% appendices are, therefore, not great, but the contents are there.

\maketitle

\noindent Dear Members of the Hiring Committee,\\[0.2cm]

I am writing to apply for the faculty position in Applied Economics at the Hong Kong University of Science and Technology. 
I am currently completing my Ph.D.\ in Economics at Bielefeld University and Università Cattolica del Sacro Cuore (EPOC Program, Marie Skłodowska-Curie Grant No.~956107). 
My research lies at the intersection of applied microeconomics, digital economics, and computational methods. 
I study how technological change, information frictions, and institutional environments shape decision-making by firms and consumers, using a combination of empirical analysis, microeconomic modeling, and machine-learning tools. 
I am particularly excited about HKUST’s strong international environment, its interdisciplinary Social Science Division, and its commitment to high-quality applied research across economics, policy, and data-intensive methods.\\

I began my Ph.D.\ in 2021, a period in which digitalization intensified across industries and generated new challenges and opportunities for applied economic analysis. 
My dissertation studies how firms and consumers respond to technological shocks, digital infrastructures, and capacity constraints in modern service markets. 
In my main chapter, I develop and estimate an agent-based model of digital service provision in which firms adjust capacity investments and pricing while heterogeneous users respond to congestion, service quality, and information flows. 
Using firm-level data and simulation-based econometrics, I estimate behavioral parameters and evaluate counterfactual technological environments to understand how digital transformation shapes productivity, market outcomes, and welfare. 
Across all three chapters—each single-authored—I combine microeconomic theory with empirical and computational methods, including Monte Carlo simulations and neural-network surrogate models. 
This framework enables the analysis of strategic behavior, learning, and adjustment dynamics in settings where standard closed-form models may be too restrictive.\\

Beyond my dissertation, my coauthored papers—\textit{Polytope Fraud Theory} and \textit{Agentic Workflows for Economic Research}—extend my interests into applied data analysis, institutional behavior, and the use of artificial intelligence in empirical economics. 
The first paper applies machine-learning techniques to detect behavioral irregularities and financial anomalies, showing how digital traces can inform oversight, governance, and applied policy questions. 
The second develops a framework for integrating large-language-model agents into empirical workflows, demonstrating how AI tools can support data collection, text analysis, and reproducible research. 
These projects reflect my broader commitment to combining economic reasoning with modern data science to address applied questions across diverse domains.\\

In future work, I plan to develop a broad applied research agenda focused on technological adoption, digital markets, and behavioral responses to information environments. 
I aim to study evolving questions related to platform competition, consumer search, market transparency, and human–AI interaction, using a combination of empirical methods, structural models, and computational tools. 
I am also interested in extending my research to applied topics in labor, public finance, and institutional economics, especially where technological change or digital data availability creates new identification possibilities. 
These directions complement HKUST’s strengths in applied microeconomics, digital economics, and interdisciplinary social science.\\

My empirical work will build on access to large administrative, financial, and digital-activity datasets, including resources available through my collaborations with Southern University of Science and Technology (SUSTech) and Bielefeld University. 
My doctoral research, funded by the \textit{Marie Skłodowska-Curie Actions} and the \textit{Bielefeld Young Researchers’ Fund}, has strengthened my ability to manage externally funded projects, work across institutions, and design computationally intensive empirical pipelines. 
I look forward to bringing these skills to HKUST’s research environment and to pursuing collaborative opportunities with faculty across the Social Science Division.\\

Thank you so much for considering my application! 
I would be honored to join the Hong Kong University of Science and Technology as a faculty member and contribute to its vibrant research and teaching community in Applied Economics. 
I look forward to the opportunity to develop an independent and impactful research program, teach at both undergraduate and graduate levels, and collaborate with colleagues across HKUST’s interdisciplinary environment.

\vspace{3mm}

\noindent Sincerely,\\
\vspace{10mm}

\noindent Zhongli Wang

\end{document}